\section{Sheaves}

\subsection{Sheaves}

\bdefn

    For a topological space $X$, its set of open subsets $\open(X)$ is partially ordered by inclusion, and thus forms a small category.
    It is also a cofiltered category: for every $U,V\in\open(X)$ they have a lower bound $U\cap V$.

\edefn

\bdefn

    Let $X$ be a topological space, and $\C$ a category.
    A {\emphcolor $\C$-valued presheaf on $X$} is a contravariant functor $\F\colon\open(X)^\op\to\C$.
    It is customary to denote $\F(U\subseteq V)\colon\F(V)\to\F(U)$ by $\res^V_U$.
    When $\C$ is a concrete category (i.e.\ its objects are sets), elements of $\F(U)$ are called {\emphcolor sections of $\F$ over $U$},
    and for a section $s\in\F(V)$ we write $s|_U$ for $\res^V_U(s)$.

    A {\emphcolor morphism of presheaves} is a natural transformation.

    Denote the category of $\C$-valued presheaves on $X$ by $\Psh_\C(X)$.

\edefn

\bexam

    Let $E$ be a set (or a topological space, or any mathematical object which lies in the same category as $X$ and its open subsets).
    We can then define the presheaf of functions on $E$ by $\F(U)=\hom(U,E)$: the set of all functions $U\to E$.
    The restriction maps are very literally restriction: for $U\subseteq V$ and $f\colon V\to E$, we restrict $f$ to $U$ to obtain a function
    $f|_U\colon U\to E$.

\eexam

\bdefn

    Let $\F$ be a presheaf over $X$ valued in a category $\C$ with filtered colimits, and $p\in X$.
    Then $\U(x)=\set{U\subseteq\hbox{ open}}[p\in U]$, the set of neighborhoods of $p$, is a cofiltered subcategory of $\open(X)$, and so $\F$
    restricts to a functor from $\U(x)^\op$.
    Define $\F$'s {\emphcolor stalk} at $p$ to be the filtered colimit
    $$ \F_p = \colim_{U\in\U(p)}\F(U) \in \C . $$

    For concrete categories, this is the set of {\emphcolor germs} of $\F$ at $p$:
    $$ \F_p = \set{(U,s)}[U\in\U(p),\ s\in\F(U)] , $$
    modulo the relation where we identify $(s,U)$ and $(t,V)$ if there is a subset $p\in W\subseteq U,V$ such that $s|_W=t|_W$.

\edefn

Taking colimits is functorial, so the correspondence $p\mapsto\F_p$ is functorial.
Explicitly $\phi\colon\F\to\G$ is a morphism of sheaves over $X$ valued in a concrete category, and $p\in X$ we define $\phi_p\colon\F_p\to\G_p$
by
$$ \phi_p(U,s) = (U,\phi_U(s)) . $$

\bdefn

    A presheaf $\F$ on $X$ (valued in a concrete category) is a {\emphcolor sheaf} if it satisfies the {\emphcolor sheaf axiom}:
    for every open subset $U\subseteq X$ and open cover $\set{U_i}$ of $U$ and sections $s_i\in\F(U_i)$, if $s_i|_{U_i\cap U_j}=s_j|_{U_i\cap U_j}$
    for all $i,j$ then there exists a {\it unique\/} $s\in\F(U)$ such that $s|_{U_i}=s_i$ for all $i$.

\edefn

We can split the sheaf axiom in two:
\benum
    \item the {\emphcolor identity axiom}: if $\set{U_i}$ is an open cover of $U$, and $s_1,s_2\in\F(U)$ are such that $s_1|_{U_i}=s_2|_{U_i}$
    for all $i$, then $s_1=s_2$.
    \item the {\emphcolor gluability axiom}: the sheaf axiom without requiring uniqueness.
\eenum

We can also state the sheaf axiom categorically.
For every open cover $\set{U_i}$ of $U$, we can consider the following diagram:
$$ \F(U) \xto\res \prod_i\F(U_i) \xtto{\res_i}[\res_j] \prod_{i,j}\F(U_i\cap U_j) , $$
where $\F(U)\xto\res\prod_i\F(U_i)$ is the map $s\mapsto(s|_{U_i})_i$, $\res_1$ and $\res_2$ are the maps
$$ \res_1(s_i)_i = (s_i|_{U_i\cap U_j})_{i,j},\qquad \res_2(s_i)_i = (s_j|_{U_i\cap U_j})_{i,j} . $$
Categorically, $\res_1$'s coordinate to $\F(U_i\cap U_j)$ composes the $i$th projection with $\res^{U_i}_{U_i\cap U_j}$, and $\res_2$'s
coordinate composes the $j$th projection with $\res^{U_j}_{U_i\cap U_j}$.
These maps are defined for any presheaf, and we can see that $\F$ is a presheaf iff this is an equalizer diagram.

Indeed, the equalizer of $\res_1,\res_2$ are all $(s_i)_i$ such that $s_i|_{U_i\cap U_j}=s_j|_{U_i\cap U_j}$, i.e.\ a family of sections which
agree on the overlaps.
And there is an $s\in\F(U)$ whose projections $s|_{U_i}$ are $s_i$ when the gluing axiom is satisfied, and it is unique (so an embedding, and thus
an equalizer) when the identity lemma is satisfied.

This allows us to generalize sheaves to any category $\C$ with arbitrary products.

The category of $\C$-valued sheaves on $X$ is denoted $\sh_\C(X)$, it is a full subcategory of $\Psh_\C(X)$; meaning sheaf morphisms are just
presheaf morphisms between sheaves.
The stalk of a sheaf at a point is just the stalk of it as a presheaf.

\bexam

    If $\F$ is a presheaf on $X$, and $U\subseteq X$ is open, we define $\F$'s {\emphcolor restriction} to $U$ to be the presheaf $\F|_U$
    on $U$ defined for $V\subseteq U$ by $\F|_U(V)=\F(V)$ (it is just the restriction of the functor $\F$ to the subcategory $\open(U)$).
    If $\F$ is a sheaf, its restriction is clearly a sheaf as well.

\eexam

\bexam

    Let $X$ be a topological space, $p\in X$, and $S$ a set (or a group, module, etc.).
    We define the {\emphcolor skyscraper sheaf} $p_*S$ by
    $$ p_*S(U) = \cases{S & if $p\in U$\cr *& else}, $$
    where $*$ is a single-point set; the restrictions are given by the identities or the unique morphism $S\to*$.
    This is indeed a sheaf: if $p\in U$ and $\set{U_i}$ is a covering of $U$ with $s_i\in\F(U_i)$ such that $s_i|_{U_i\cap U_j}=s_j|_{U_i\cap U_j}$
    then for every $p\in U_i$ we have a single value $s\in S$ which is equal to all $s_i\in\F(U_i)$, this is the glued section.
    Otherwise $p\notin U$ and $s_i$ are all the unique element in $*$, so $s$ must be too.

\eexam

\bexam

    Let $X$ be a topological space and $S$ a set.
    The {\emphcolor constant presheaf} associated with $S$ is the presheaf $\ul S_\pre$ which maps all open sets $U$ to $S$.
    This is a presheaf, but not a sheaf in general.
    For example if $X=\set{p,q}$ is the two point space with the discrete topology then take $s_p\neq s_q$; they cannot be glued to a global section
    on $X$.

    If we instead define $\ul S(U)$ to be the set of locally constant functions $U\to S$, then this is a sheaf, called the {\emphcolor constant
    sheaf}.
    It is just the sheaf of continuous functions $\ul S(U)=C(U,S)$ with $S$ given the discrete topology.
    The sheaf of continuous functions $C(\bul,Y)$ is a sheaf in general (by the gluing lemma from topology).

\eexam

\bexam

    Let $\pi\colon X\to Y$ be a continuous map, and $\F$ a presheaf over $X$.
    Define the {\emphcolor pushforward} (or {\emphcolor direct image}) presheaf $\pi_*\F$ by
    $$ \pi_*\F(V) = \F(\pi^{-1}V), $$
    for $V\subseteq Y$, and $\res^V_W=\res^{\pi^{-1}V}_{\pi^{-1}W}$.
    This is clearly a presheaf; if $\F$ is a sheaf it is also a sheaf.
    Indeed, if $V\subseteq Y$ is open and $\set{V_i}_i$ is a cover with $s_i\in\pi_*\F(V_i)=\F(\pi^{-1}V_i)$.
    Then the glue of these is unique, because $\F$ is a sheaf.

    The skyscraper sheaf $p_*S$ is the pushforward of the constant sheaf $\ul S$ by the inclusion $\set p\to X$.

    An important property of the direct image sheaf is that it is functorial: if $\phi\colon\F\to\G$ is a morphism of presheaves over $X$,
    then $\pi_*\phi\colon\pi_*\F\to\pi_*\G$ defined by $\pi_*\phi_V=\phi_{\pi^{-1}V}$ is a morphism of presheaves.

\eexam

Note that if $\pi\colon X\to Y$ is continuous, $\F$ is a sheaf over $X$, and $\pi(p)=q$, then there is a natural morphism of stalks
$$ (\pi_*\F)_q \to \pi_p . $$
Indeed, the left-hand side is
$$ \colim_{q\in V}\F(\pi^{-1}V) = \colim_{p\in\pi^{-1}V}\F(\pi^{-1}V) . $$
Now a cone under $\F$ defines a unique cone under $\pi_*\F$ simply by restricting, and since $\colim\pi_*\F$ is the initial cone under $\pi_*\F$
there is a canonical map to this restricted cone, in particular for $\colim\F$.

In more concrete terms, $(\pi_*\F)_q$ has elements $(s,V)$ for $s\in\F(\pi^{-1}V)$, which we then map to $(s,\pi^{-1}V)\in\F_p$.

\bdefn

    If $\O_X$ is a sheaf of rings over a topological space $X$ (i.e.\ a sheaf over $X$ valued in the category of rings), then the pair
    $(X,\O_X)$ is called a {\emphcolor ringed space}; $\O_X$ is called the {\emphcolor structure sheaf}.
    Sections of $\O_X$ over $U\subseteq X$ are called {\emphcolor functions} on $U$, global sections are caled {\emphcolor global functions}.

\edefn

\bdefn

    If $\O_X$ is a ringed space, an {\emphcolor $\O_X$-module} is a sheaf of Abelian groups $\F$ over $X$ such that for each $U\subseteq X$,
    $\F(U)$ is an $\O_X(U)$-module, and all of these modules are well-behaved in that when $U\subseteq V$, the following diagram commutes:

    \commsq{\O_X(V)\times\F(V)&\F(V)\cr \O_X(U)\times\F(U)&\F(U)}
        {{\rm mult}}{{\rm mult}}{\res_{V,U}\times\res_{V,U}}{\res_{V,U}}

    The horizontal arrows are $\O_X(U)$'s action on $\F(U)$.

\edefn

The commutative diagram means that
$$ (s\cdot x)|_U = (s|_U)\cdot(x|_U) , $$
for $s\in\O_X(V),x\in\F(V)$.

\bprop

    If $\F$ is an $\O_X$-module, then for each $p\in X$, $\F_p$ is an $\O_{X,p}$-module.

\eprop

\bproof

    The condition that $\F$ be an $\O_X$-module means that multiplication is a natural map $\O_X\times\F\to\F$.
    Colimits are taken pointwise in functor categories (in particular for the product of functors), so this gives a natural map
    $\O_{X,p}\times\F_p\to\F_p$.
    The natural map is given by mapping $(s,U),(x,V)$ to $(sx,U\cap V)$, which is easily seen to give a module structure.
    \qed

\eproof

\bexam

    Consider a smooth manifold $X$; this defines a ringed space $(X,\O_X)$ where $\O_X$ is the sheaf of smooth real-valued functions:
    $\O_X(U)=C^\infty(U)$.
    These are all real algebras, so we could view this as an $\bR$-ringed space.

    If $\pi\colon V\to X$ is a vector bundle over $X$, then consider the sheaf of smooth local sections of $\pi$:
    $$ \Gamma(U) = \set{\sigma\colon U\to V\hbox{ smooth}}[\pi\circ\sigma=\id_U] . $$
    Then $\Gamma$ is an $\O_X$-module: it is a classical result that a local section defined on $U$ can be multiplied by a smooth real-valued
    function defined on $U$.

\eexam

\bnote

    Instead of $\sh_\C(X)$ to denote the category of $\C$-valued sheaves on $X$, we will denote the category by $\C_X$.
    And the category of presheaves by $\C_X^\pre$.

\enote

Note that if $\F,G\in\C_X$ are sheaves (or presheaves), then the hom-set of sheaf morphisms $\F\to\G$ is indeed a set, since $X$ as a category
is small (functor categories under small categories are locally small).
Furthermore, if $\C$ is enriched over some monoidal category (in our case, usually $\Ab$) so is $\C_X$.
This is because it is true in general for functor categories over enriched categories.

\bdefn

    Let $\F,\G$ be $\C$-valued sheaves defined on $X$.
    Define their {\emphcolor sheaf-hom} to be the sheaf whose sections are sheaf morphisms,
    $$ \shom(\F,\G)(U) = \hom_{\C_X}(\F|_U,\G|_U) . $$
    As discussed, if $\C$ is enriched over a monoidal category $V$ then so is $\C_X$, and so $\shom(\F,\G)$ takes values in $V$.
    In particular, if $\F,\G$ are sheaves of Abelian groups, then so is their sheaf-hom.

\edefn

\bprop

    The sheaf-hom is indeed a sheaf.

\eprop

\bproof

    We first show it is a sheaf: if $U\subseteq V$ then a sheaf morphism $\F|_V\to\G|_V$ restricts to a sheaf morphism
    $\F|_U\to\G|_U$ in the obvious way: its components are the same but only taken on subsets of $U$.
    This is clearly a presheaf.

    To show that it is a sheaf, let $\set{U_i}$ be an open cover for $U\susbeteq X$, and $\alpha_i\colon\F|_{U_i}\to\G|_{U_i}$ be sheaf morphisms
    such that $\alpha_i|_{U_i\cap U_j}=\alpha_j|_{U_i\cap U_j}$ for all $i,j$.
    In order to glue them together for a sheaf morphism $\alpha\colon\F_U\to\G_U$ we need to define $\alpha_W$ for all $W\subseteq U$.
    For $s\in\F(W)$ consider $s_i=s|_{U_i}\in U_i\cap W$, and its image under $\alpha_i$ to $\G(U_i\cap W)$, call it $t_i$.
    Then since $U_i\cap W$ cover $W$ and since $\alpha$ is a morphism as well as $\G$ being a sheaf, we can uniquely glue $t_i$ together to
    get a $t\in\G(W)$; this will be the image of $s$.

    It is not hard to see that this is the unique gluing of $\alpha_i$, and is itself a sheaf morphism.
    \qed

\eproof

\bnote

    In the above proof we needed only for $\G$ to be a sheaf, $\F$ can be a presheaf.

    Furthermore $\shom$ forms a bifunctor
    $$ \shom\colon\C_X^\op\times\C_X \to \Set_X , $$
    or $V_X$ instead of $\Set_X$ when $\C$ is enriched over $V$.

\enote

Note that if $\O_X$ is a ringed space, it is also a module over itself.
Thus if $\F$ is an $\O_X$ module, we can form its {\emphcolor dual}:
$$ \F^\vee = \shom_{\Mod_{\O_X}}(\F,\O_X) . $$

\subsection{Sheafification}

There is a natural map
$$ \F(U) \to \F_p , $$
for every $p\in U$ given by mapping $s\in\F(U)$ to $(U,s)\in\F_p$.
Together these assimilate into a map
$$ \F(U) \to \prod_{p\in U}\F_p . $$

\bprop

    The natural map above is injective when $\F$ is a sheaf.

\eprop

\bproof

    Suppose that $s,t\in\F(U)$ such that $(U,s)=(U,t)$ in $\F_p$.
    Then there exists a neighborhood $p\in W_p\subseteq U$ such that $s|_{W_p}=t|_{W_p}$.
    If this is true for all $p$, then $W_p$ form an open cover of $U$ and by the identity axiom, $s=t$.
    \qed

\eproof

\bdefn

    An element $(s_p)_p\in\prod_{p\in U}\F_p$ consists of {\emphcolor compatible germs} if for all $p\in U$, $s_p$ has a representative
    $\tilde s_p\in\F(U_p)$ for $p\in U_p\subseteq U$ such that the germ of $\tilde s_p$ at $q\in U_p$ is $s_q$ for all $q\in U_p$.

\edefn

Equivalently, this says that there is an open cover $\set{U_i}$ of $U$ and sections $f_i\in\F(U_i)$ such that if $p\in U_i$ then $s_p$ is the germ
of $f_i$ at $p$.
If $s$ is a section of $\F$ over $U$, then $(s)_p$ consists of compatible germs.

\bprop

    The natural map $\F(U)\to\prod_{p\in U}\F_p$'s image is precisely the compatible germs.

\eprop

\bproof

    We already showed that every tuple in its image consists of compatible germs, now we show the converse.
    Suppose $(s_p)_p$ consists of compatible germs, so let $\tilde s_p\in\F(U_p)$ be a representation as in the definition.
    Then $U_p$ cover $U$, and $\tilde s_p|_{U_p\cap U_q}$'s germ at $q$ is $s_q$ so that $\tilde s_p,\tilde s_q$ agree on a neighborhood of $q$
    and by symmetry in a neighborhood of $p$ and $q$.

\eproof

